\section{Finite iteration and exact boundary maps}
\label{iter:section-boundaries}

We work with schemes whose arrows are componentwise strong system
embeddings, preserving the specified distinguished data. All levels in
this section are eligible strong limit cardinals of countable cofinality.
The additional closure requirements below are imposed for each fixed
finite iteration depth. A common choice for all finite depths will be
made separately in Theorem~\ref{close:common-levels}.
Unboundedness after the refinements, including the supply of smaller
completion levels for exact lifting, uses the eligible-subclass
hypothesis of Convention~\ref{lev:eligible-subclass}.  In the finite
marked application, Section~\ref{close:section-uniform} verifies this
hypothesis by maintaining containment of $G_j$ for some finite $j$, in
the hierarchy previewed in Section~\ref{lev:section-working-schemes}.

Completion and uniqueness for finite systems are central in
multidimensional independence.  Shelah developed dimension-raising
existence and uniqueness methods in \cite[Sections~3--5]{Shelah1983B};
compare the distinction between extension, uniqueness, and strong
uniqueness in \cite[Definition~8.16]{ShelahVasey2024}. Here amalgamated completion
and extension onto the prescribed target are proved separately.

\subsection{Forgetting components and recovering types}

Iteration requires more than forgetting vertices of an object: a type
prescribed in the root must determine a unique type over the entire
system.  The following correspondence supplies that requirement and
shows that it is compatible with free lifting.

\begin{proposition}[Type correspondence under projection]
\label{iter:projection-types}
Let $P:\mathcal C\longrightarrow\mathcal D$ forget some components of
systems. Suppose that both schemes satisfy the working package,
that $P$ takes objects and arrows to objects and arrows, including
mixed arrows, and that the following diagram comparison holds: for
bounded tuples in the retained sorts, equality of their
$\mathcal C$-diagrams is equivalent to equality of their
$\mathcal D$-diagrams after projection. The comparisons are made at
arities bounded at both levels when levels differ.

At sufficiently closed levels, restriction induces a bijection
on types of bounded tuples in any fixed list of retained sorts,
given by
\[
 \tp_{\mathcal C}(v/X)\longmapsto\tp_{\mathcal D}(v/PX).
\]
The parameters on the left are all of $X$, including its forgotten
components. The bijection commutes with restriction and isomorphisms,
and with free lifting along arrows, including mixed arrows when the
tuple arities and the required testing bounds lie below the smaller
level. It also takes witnessed free products to their projected
witnessed free products.
\end{proposition}

\begin{proof}
First consider surjectivity at a level $\mu$. Present
$X=F(I)$ on a well-order, and let $p$ be a projected type over $PX$.
For a bounded list $a$ of parameters from $PX$, saturation in
$\mathcal D$ realizes the prescribed joint diagram of $(v,a)$ in
$PX$ itself. Such a realization has term representations in the
retained sorts of $F(I)$. Thus every one of the simultaneous local
requests used in compact placement has an actual placement in the
chart of $X$.

Here we prescribe projected diagrams only; diagrams involving the
forgotten sorts have yet to be chosen. The proof of
Lemma~\ref{calc:compact-placement} applies to this form of request as
well. To verify this explicitly, fix a sufficiently large bounded
cardinal $\chi=\chi^{|v|+\aleph_0}$, larger than the label and testing
bounds. A profile records the labels of $v$, its ordered supports,
and their comparisons with the chosen representations of the
projected parameters under consideration. There are at most $\chi$
profiles over at most $\chi$ parameters. If a continuation of a
profile fails a prescribed projected diagram, record a witnessing
projected test and add it to the next parameter set. The comparison
of continuations with no new event uses equality of natural
$\mathcal C$-diagrams on retained sorts; the hypothesis on $P$
then gives equality of the requested $\mathcal D$-diagrams.
Consequently the same new-hit or new-separation argument as in
Lemma~\ref{calc:compact-placement} rules out failure at every step.
The local possibility used at the final step is the simultaneous
request on the entire final bounded parameter set, which is
realizable in $PX$ by saturation.

The resulting placement in an enlargement of $I$ satisfies the
projected requests up to $\chi$. Taking $\chi$ beyond the testing
bound in $\mathcal D$ makes its projected type equal to $p$.
The placement itself defines a $\mathcal C$-type over all of $X$,
which proves surjectivity. In particular, no diagram on a forgotten
parameter was assumed to have been prescribed in advance.

Next let $q=\tp_{\mathcal C}(v/X)$. Its approximation lists from
Lemma~\ref{calc:approximations} consist of tuples in the same
retained sorts as $v$. A query whose parameters lie in the projection
sorts can be moved over such an approximation list using the
working package. The projection comparison identifies precisely
the equality tests used in this computation. The majority answer
and its exception estimate therefore compute the projected free
lift. Increasing the approximation length or the test arity gives
the same answers by the compatibility part of the approximation
construction. Thus free lifts project to free lifts. This reasoning
uses bounds below the originating level and consequently applies
also to mixed arrows. Projecting the witnessing chain of a free
product now gives its claimed projected witnessing chain.

For injectivity, suppose $q_0\ne q_1$ have the same projection.
Choose a bounded tuple $f$ from $X$ witnessing their inequality:
the diagrams of $(v,f)$ for the two types have distinct values
$\delta_0$ and $\delta_1$. Choose a bounded infinite $\chi$ such
that
\[
 \chi^{|f|+\aleph_0}=\chi,
\]
and such that $\chi$ bounds $|v|$ and the label and orbit-counting
costs. All constructions below fit below $\mu$: in particular
$2^\chi<\mu$, and a tuple of $\chi$ many copies of $v$ remains
bounded at $\mu$.

For each $\eta\in{}^\chi 2$, construct a genuine witnessed free
sequence $(v_i^\eta:i<\chi)$ over $X$ with successive marginal
$q_{\eta(i)}$. Successor steps use free extension and amalgamation;
limit steps use the chain axiom. The uniqueness of whole types of
witnessed chains in Lemma~\ref{iso:witnessed-uniqueness} applies also
to these bounded tuple entries. Their projected marginal types
are all the same, so their projected sequences have the same whole
type over $PX$. In particular their parameter-free projected
diagrams agree. By the hypothesis on $P$, their whole
$\mathcal C$-diagrams agree as well.

Fix one such sequence $a=(a_i:i<\chi)$ in a level-$\mu$ object.
Bounded diagram homogeneity gives, for every $\eta$, an isomorphism
from the object containing $v^\eta$ onto this fixed object, carrying
$v_i^\eta$ to $a_i$ for every $i$. Write $f_\eta$ for the image
of $f$. Preservation of diagrams gives
\[
 D_{\mathcal C}(a_i,f_\eta)=\delta_{\eta(i)}.
\]
If $\eta\ne\eta'$, a coordinate where they differ shows that
$f_\eta$ and $f_{\eta'}$ have different orbits over the entries
of $a$. There would be $2^\chi$ such orbits of $|f|$-tuples over
at most $\chi$ parameters. Lemma~\ref{calc:parameter-orbits}
bounds their number by $\chi$, a contradiction.

Restriction and transport commute with projection directly.
Surjectivity and injectivity, together with the free-lift calculation,
give all the stated compatibilities. In particular a projected
extension of a type is free exactly when its corresponding full
extension is free: the free full extension has the required
projection, and injectivity makes it the only possible full lift.
\end{proof}

\begin{corollary}[Lifting an extension of a projection]
\label{iter:projection-extension}
Under the hypotheses of Proposition~\ref{iter:projection-types},
let $X$ and $Z$ be at the same eligible level, and let
$PX\longrightarrow Z$ be an arrow. There are a
$\mathcal C$-extension $X\longrightarrow X^*$ and a
$\mathcal D$-arrow $Z\longrightarrow PX^*$ whose composite on
$PX$ is the projection of $X\longrightarrow X^*$.
\end{corollary}

\begin{proof}
Identify the given base with its image. Choose nested bounded
lists $d_i$ exhausting $Z$ for $i<\omega$. In each vertex include
enumerations of increasing smaller strong submodels exhausting
that vertex. We construct extensions $X\to X_i$ and assignments
of $d_i$ into $PX_i$ that extend the previous assignments and
have the correct type over $PX$.

Suppose the assignment of $d_i$ is given. Amalgamation in
$\mathcal D$, with identification of the already matching tuple,
puts $Z$ and $PX_i$ in a common extension, agreeing on $PX$ and
on $d_i$. The type there of $d_{i+1}$ over $PX_i$ lifts by
Proposition~\ref{iter:projection-types} to a type over $X_i$.
Realize it in a $\mathcal C$-extension $X_{i+1}$. The part
corresponding to $d_i$ is already a tuple of parameters from
$PX_i$, so its equality diagram forces the new assignment to
extend the old one. At the first step use the same argument with
the empty tuple.

Take the union of the $X_i$. The assigned lists match all their
bounded diagrams, and the included vertex submodel enumerations
give componentwise strongness by the working-package map
criterion. Thus their union is the required arrow from $Z$.
\end{proof}

\subsection{The finite systems and their faces}

\begin{theorem}[Finite marked systems]
\label{iter:finite-systems}
Start with the single-model scheme, allowing fixed bounded named
data of a fixed diagram. Let $r$ be a prescribed, definable,
isomorphism-invariant stationary type of a finite tuple in its
root sort. Stationarity here includes free lifting under mixed
arrows at eligible levels. For each finite $n$, the $n$-fold marked
increment construction gives a working scheme $\mathcal C_n$
at sufficiently closed levels. Its objects are systems indexed
by the subsets of $\{1,\ldots,n\}$, with the prescribed marks
and componentwise strong arrows.

For $s\subseteq\{1,\ldots,n\}$, retaining the axes in $s$
in their induced order gives a projection
\[
 P_s:\mathcal C_n\longrightarrow\mathcal C_{|s|}.
\]
Every such projection takes objects and arrows to objects and
arrows, including mixed arrows, and preserves and reflects
diagram equality on bounded tuples in its retained sorts.
It consequently has the exact type correspondence of
Proposition~\ref{iter:projection-types}. The type in the root
sort lying over $r$ exists uniquely over every whole system and
lifts freely.

These conclusions include the projection to the base root and
each single-axis projection. In a single-axis projection the
upper object is all-dominated over its lower object together
with the prescribed mark; forgetting other axes does not remove
that assertion.
\end{theorem}

\begin{proof}
We induct on the number of axes. At depth zero the working
package and the prescription $r$ are the given ones. Suppose
the assertion is known for $\mathcal C_n$. By the exact type
correspondence for the root projection, $r$ determines a unique
type over every $\mathcal C_n$-object. Projection commutes with
free lifting, so this prescription is stationary, including
under mixed arrows. Theorem~\ref{mark:package} therefore supplies
the working scheme
$\mathcal C_{n+1}=(\mathcal C_n)^b$.

First forget its last axis. The result is its lower object
$X$. If two bounded tuples in the lower objects have equal
$\mathcal C_n$-diagrams, a lower-object isomorphism matches them.
It transports the prescribed type of $b$, and hence extends
to an isomorphism of the marked pairs by
Lemma~\ref{mark:pair-isomorphism}. This proves diagram reflection;
preservation follows by restriction. Objects and arrows plainly
project to their lower objects and arrows.

Next apply an earlier projection
$P:\mathcal C_n\to\mathcal C_k$ to both layers of
$(X\le Y,b)$. The marginal of $b$ over $PX$ is the required
type, by Proposition~\ref{iter:projection-types}. We must verify
all-domination of $PY$ over $PX+b$. Suppose it fails for a
bounded tuple $a$ from $PY$. The nonisolation alternative of
Proposition~\ref{iso:alternative} gives
an extension of
\[
 \tp_{\mathcal C_k}(b,a/PX)
\]
to an enlarged parameter object $V\ge PX$ that is nonfree,
although its $b$-marginal is free. By
Corollary~\ref{iter:projection-extension}, put $V$ inside the
projection of an extension $X'$ of $X$. Freely extend the
chosen joint type to $PX'$. It remains nonfree from $PX$,
since otherwise its restriction to $V$ would have been free.
Its $b$-marginal remains free from $PX$ by transitivity.

Lift this type to $X'$ by the exact type correspondence. Its
restriction to $X$ is the original type of $(b,a)$, because
their projections agree and the correspondence is injective.
Its $b$-marginal is the free lift of the original one, again
by injectivity and compatibility with free lifting. Its joint
type is nonfree: a free joint type would project to a free
one. This contradicts the relative isolation of
$\tp_{\mathcal C_n}(b,a/X)$ and
Lemma~\ref{iso:free-isolation}. The projected pair is thus
all-dominated as required. Componentwise projection of an arrow
preserves the prescribed mark and gives an arrow of these pairs.

It remains to reflect pair diagrams. Suppose bounded retained
tuples in two projected pairs have equal diagrams. The projected
pair isomorphism test gives an isomorphism of their projected
bases transporting their types jointly with $b$; any indicated
base entries are included in this match. Choose in the first
base sufficiently long approximation lists for this joint type,
at arities exceeding the full and projected testing bounds.
Their entries are in retained sorts. Carry these lists, together
with the indicated base entries, by the projected base
isomorphism. The inductive reflection hypothesis and bounded
homogeneity give an isomorphism of the full bases making
precisely these bounded matches.

The transported approximation lists compute the same projected
joint type over the second projected base. Exact type
correspondence therefore says that the full base isomorphism
transports the entire joint type with $b$. The full pair
isomorphism test now extends it to an isomorphism of marked
pairs matching the original tuple. Preservation is immediate
from projection of an isomorphism. The same calculation uses
free lifting and the mixed preservation and reflection in
Lemma~\ref{mark:pair-isomorphism} when comparing eligible levels;
all the tuple and approximation bounds are below the smaller
one. Composing elementary axis deletions proves the assertion
for every $P_s$. This completes the finite induction.
\end{proof}

\begin{definition}[Boundary]
\label{iter:boundary}
For an $n$-axis object $J$, its boundary $\partial J$ consists
of all the proper-axis projections $P_sJ$, with their specified
inclusions. For $n=0$ the boundary is empty. A boundary map
$h:\partial J\to\partial K$ is a compatible family of arrows
$h_s:P_sJ\to P_sK$ for proper $s$. A boundary isomorphism is
such a family of isomorphisms.

Lists in a boundary are lists of occurrences in its vertex
sorts. The diagram of a joint list records the specified
inclusion maps and all resulting equalities. In particular,
entries supplied from different faces can coincide in a
containing vertex. Compatibility on the named subfaces alone
is not asserted to settle these extra equalities; their
preservation will be part of the next argument.
\end{definition}

\subsection{Completing a boundary into an extension}

\begin{lemma}[Joint diagrams on the boundary]
\label{iter:boundary-diagrams}
At the levels used for $\mathcal C_n$, every boundary map
between two $n$-axis objects at the same level preserves all
bounded $\mathcal C_n$-diagrams on their boundary unions.
This includes joint diagrams of entries from different faces
and all equality information on their actual overlaps.
\end{lemma}

\begin{theorem}[Amalgamated boundary completion]
\label{iter:boundary-completion}
For each finite $n$, at sufficiently closed levels the following
holds. If $J$ and $K$ are $\mathcal C_n$-objects at the same
level and $h:\partial J\to\partial K$ is a boundary map,
there are arrows
\[
 K\longrightarrow K^*,\qquad J\longrightarrow K^*
\]
in $\mathcal C_n$ such that the second restricts to the
composite of $h$ with the first on every proper projection.
The extension is in the same scheme and at the same level.
No assertion about completion between distinct levels is
needed here.
\end{theorem}

\begin{proof}[Proof of Lemma~\ref{iter:boundary-diagrams} and
Theorem~\ref{iter:boundary-completion}]
We prove the two assertions together by induction on $n$.
For $n=0$ the diagram assertion is empty, and joint embedding
gives completion. At a positive dimension, after proving the
diagram assertion we shall derive completion by a placement
argument. Thus only lower-dimensional completion will be used
in proving the next diagram assertion.

For $n=1$, its only proper projection is its base.
Theorem~\ref{iter:finite-systems} identifies the full-scheme
diagrams on that base with its projected diagrams. The given
base arrow preserves the latter, proving the assertion.

Let $n=m+1$ with $m\ge1$, and write the source and target as
\[
 (X^-\le Y^-,b^-),\qquad (X^+\le Y^+,b^+)
\]
over $\mathcal C_m$. A boundary map gives a full base arrow
$e:X^-\to X^+$ and a boundary map
$h:\partial Y^-\to\partial Y^+$ agreeing with $e$ on
$\partial X^-$; the upper faces also specify
$h(b^-)=b^+$. Let $v$ be a bounded tuple from
$\partial Y^-$, enlarged to include $b^-$.
We first claim that its type over $X^-$ in $\mathcal C_m$
equals the restriction along $e$ of the type of $h(v)$.

Suppose instead that a bounded tuple $f$ from $X^-$ witnesses
a difference. Put
\[
 \delta_0=D_{\mathcal C_m}(f,v),\qquad
 \delta_1=D_{\mathcal C_m}(e(f),h(v)),
 \qquad \delta_0\ne\delta_1.
\]
Choose an infinite bounded $\chi$ satisfying
$\chi^{|f|+\aleph_0}=\chi$, bounding $|v|$ and all the
label and orbit-counting costs. Let $\epsilon$ be the least
cardinal with $2^\epsilon>\chi$. Cantor's Theorem gives
$\epsilon\le\chi$, and minimality gives $2^{|i|}\le\chi$
for every $i<\epsilon$. Consequently
\[
 \left|{}^{<\epsilon}2\right|\le\chi,
 \qquad 2^\epsilon\le 2^\chi<\mu,
\]
where $\mu$ is the common level. Every accumulator and branch
world below will have size $\mu$, and the total number of
amalgamations remains below $\mu$.

We construct branch worlds $W_s\in\mathcal C_m$ for
$s\in{}^{\le\epsilon}2$, carrying tuples $f_s$, with
coherent full arrows along each branch. Their proper
boundaries have compatible arrows
\[
 q_s:\partial W_s\longrightarrow\partial A
\]
into an increasing common accumulator $A\in\mathcal C_m$.
The accumulator may be enlarged at each operation; existing
boundary arrows are always composed with its enlargement
arrow. The invariant requires coherence of these boundary
arrows with the branch arrows. It does not require full maps
of all the branch worlds into the accumulator to be retained.

Start with $W_\varnothing=X^-$, its tuple $f$, and the
identity boundary map into $A=X^-$. At a node $s$, diagram
preservation along the branch shows that
$D(f_s)=D(f)$. Bounded homogeneity supplies an isomorphism
\[
 k_s:X^-\overset{\cong}{\longrightarrow}W_s,
 \qquad k_s(f)=f_s.
\]
There is no coherence requirement on the choices of $k_s$.
Choose the minus child to be a copy of $Y^-$ attached over
$W_s$ by $k_s^{-1}$, and the plus child to be a copy of
$Y^+$ attached by $e k_s^{-1}$. These choices give the full
branch arrows. In each child, $f$ denotes the image of $f_s$;
the corresponding upper boundary lists have diagrams
$\delta_0$ and $\delta_1$ with it. Transporting $h$ gives
a boundary arrow
\[
 h_s:\partial W_{s0}\longrightarrow\partial W_{s1}
\]
over $\partial W_s$, matching the two copies of $v$.

Apply the already proved completion theorem in dimension $m$
to $q_s$. After enlarging the accumulator, this gives a full
map $W_s\to A_0$ extending $q_s$. Ordinary amalgamation in
$\mathcal C_m$ then puts the plus child in an enlargement
$A_1$ over this full map. Let $G_1:W_{s1}\to A_1$ be the
resulting arrow. Retain its boundary restriction as $q_{s1}$,
and define
\[
 q_{s0}=q_{s1}h_s.
\]
Both restrict to the preceding $q_s$, after enlargement, and
the two boundary copies of $v$ have one common image $v_s$
in $A_1$. The temporary full maps to $A_0$ and $A_1$ are
no longer part of the invariant. In particular the minus
child has not been required to send its copy of $f$ to the
temporary image of $f_s$ in the accumulator.

\begin{figure}[htbp]
\centering
\begin{tikzpicture}[>=Latex, every node/.style={font=\small},
  world/.style={draw,rounded corners,inner sep=5pt},
  boundary/.style={draw,inner sep=4pt}]
 \node[world] (w) at (0,2.4) {$W_s,f_s$};
 \node[world] (w0) at (-2,1.1) {$W_{s0},f_{s0}$};
 \node[world] (w1) at (2,1.1) {$W_{s1},f_{s1}$};
 \node[boundary] (b0) at (-2,-.1) {$\partial W_{s0}$};
 \node[boundary] (b1) at (2,-.1) {$\partial W_{s1}$};
 \node[world] (a) at (0,-1.7) {$A_1\ni v_s$};
 \draw[->] (w) -- (w0);
 \draw[->] (w) -- (w1);
 \draw (w0) -- node[left] {boundary} (b0);
 \draw (w1) -- node[right] {boundary} (b1);
 \draw[->] (b0) -- node[above] {$h_s$} (b1);
 \draw[->,dashed] (b0) -- node[left] {$q_{s0}$} (a);
 \draw[->,dashed] (b1) -- node[right] {$q_{s1}$} (a);
 \draw[->,dotted,rounded corners=6pt] (w1.east) -- (4.2,1.1)
   -- node[right] {$G_1$} (4.2,-1.7) -- (a.east);
\end{tikzpicture}
\caption{A return-tree split. Solid arrows between worlds are
retained along their separate branches. Dashed boundary maps
persist in the common accumulator and agree on the test $v_s$.
Each dashed arrow denotes a compatible family of maps into the
proper projections of $A_1$.
The dotted full map is used to choose the plus boundary map
and need not be retained at later splits. Full maps of all
leaf worlds are installed only after the splitting tests have
been fixed.}
\label{iter:fig-return-tree}
\end{figure}

Process the nodes of each level successively. At a limit
height take the union of the worlds along each branch and
the union of the accumulators constructed so far. Their
sizes are still $\mu$. The boundary arrows into the union
accumulator exist by the chain and common-upper-bound clauses
in every proper projection. Full branch arrows and the
carried tuples also pass to the unions. This establishes
the invariant through height $\epsilon$.

At that height, apply lower-dimensional completion to each
leaf boundary map, successively enlarging the accumulator.
This supplies a full map of each $W_\eta$, for
$\eta\in{}^\epsilon2$, extending its retained boundary
map. There are fewer than $\mu$ leaves, so the final union
is still at level $\mu$. If two leaves first differ after
a node $s$, their images of $f$ have diagrams $\delta_0$
and $\delta_1$ over the common tuple $v_s$. Indeed the
corresponding copy of $v$ was in that child's boundary,
remained there along that branch, and its boundary image was never
changed except by composition with an accumulator arrow.
Thus the final images of $f$ occupy more than $\chi$
orbits over the entries of $(v_s:s\in{}^{<\epsilon}2)$.
This parameter list has size at most $\chi$. The orbit
bound of Lemma~\ref{calc:parameter-orbits} gives the
contradiction. The claimed equality of types of $v$ follows.

To pass from this type equality to the full pair diagrams, retain
an arbitrary bounded tuple $a$ from the lower face $X^-$.
The marginal of $b^+$ over
$X^+$ is the prescribed free lift from $e(X^-)$.
Relative isolation of the source type of $(b^-,v)$ and
Lemma~\ref{iso:free-isolation} therefore identify the actual type
of $(b^+,h(v))$ over $X^+$ with the free lift of that source type.

Choose sufficiently long bounded approximation lists in $X^-$
computing the source type.  Their images under $e$ compute its
free lift over $X^+$ by Theorem~\ref{calc:free-extension}.
Base homogeneity gives an isomorphism $X^-\cong X^+$ matching
$a$ and these lists with their images under $e$, since $e$
preserves their joint bounded diagram.  The transport assertion
for the approximation lists says that this isomorphism carries
the whole source type to the type of $(b^+,h(v))$.
Lemma~\ref{mark:pair-isomorphism} now extends it to an isomorphism
of the marked pairs matching $(a,v)$ with $(e(a),h(v))$.
Every bounded boundary list is of this form.  Thus its full
$\mathcal C_n$-diagram is preserved, including cross-face equalities.

We finish the induction by deriving completion from this
diagram equality. Present the source $J$ on a well-order.
Its boundary union is a natural collection of the
designated term sorts in that presentation. Choose nested
bounded lists $d_i$ exhausting $J$, including increasing
strong submodel enumerations in every vertex. We construct
extensions $K_i$ of $K$ and compatible assignments of
$d_i$ into $K_i$ satisfying
\[
 D_J(d_i,u)=D_{K_i}(d_i',h(u))
 \quad\text{for every bounded boundary list }u.
 \tag{C}
\]
Here and subsequently $h$ is composed with the current
extension arrow. At the start the assertion for the empty
list is exactly the boundary diagram equality just proved.

For the next list prescribe the source joint diagrams on
the parameter set consisting of the boundary images and
the previous assignment. Each bounded local request,
including the whole previous assigned list, is possible
by saturation: the parameter diagram agrees by $(\mathrm C)$.
This also checks repetitions and equalities between old
entries and boundary entries, so the requests define one
coherent prescription. Apply compact placement in a
well-ordered chart of $K_i$, keeping these previous
entries fixed. Work at one bounded cardinal large enough
for the profile count, the new tuple, and the testing
bound for the source boundary layout with the previous
tuple included. The resulting placement lies in a
level-$\mu$ extension $K_{i+1}$.

To check $(\mathrm C)$ beyond that initial testing cardinal,
pull the candidate diagrams back along $h$ to the natural
boundary layout of $J$. Their local possibility there
follows from the boundary diagram equality and saturation.
Lemma~\ref{calc:compact-placement} compactifies these
requests over this source layout, and
Proposition~\ref{calc:testing} says that agreement at the chosen
testing bound determines all bounded comparisons.
The actual images of the boundary inside $K_i$ have not
been assumed to form a natural layout. All uses of that
hypothesis are on the source layout after pullback.

Continue through $\omega$ steps and take the union of the
$K_i$. The working-package criterion, with the included
strong submodel enumerations, makes the union assignment
an arrow $J\to K^*$. Whenever a boundary entry is first
included in $d_i$, its equality test against that boundary
entry forces its assigned image to be $h$ of that entry.
The union arrow therefore extends the entire boundary
map. The map from $K$ is the chain inclusion. This proves
completion and closes the induction.
\end{proof}

\subsection{A bounded support for boundary tests}

We have extended boundary maps into common extensions.  Exact lifting
must instead land onto the given target.  The next lemma supplies bounded
data that control a tuple's boundary comparisons at a fixed bounded test
arity.  Together with uniform testing, this permits back-and-forth to
remain inside the two given systems.

\begin{lemma}[Fixed parameters for a boundary test function]
\label{iter:boundary-fixed-parameters}
Fix $n<\omega$. At further sufficiently closed levels,
let $J\in\mathcal C_n$, let $z$ be a bounded tuple of
$J$, and let $\sigma$ be a bounded infinite test arity.
There is a bounded list $a$ in $\partial J$ such that
every boundary automorphism $\alpha$ fixing $a$ pointwise
satisfies
\[
 D_J(z,u)=D_J(z,\alpha(u))
 \quad\text{for every boundary tuple }u
 \text{ of length at most }\sigma.
 \tag{F}
\]
A boundary automorphism means a compatible family of
isomorphisms of all proper projections; a full
automorphism of $J$ is not assumed.
\end{lemma}

\begin{proof}
The assertion is immediate for the empty boundary.
Suppose it fails otherwise. Choose a bounded infinite
$\chi$ with $\chi^{|z|+\aleph_0}=\chi$, bounding
$\sigma$ and the orbit-counting and label costs. Let
$\epsilon$ be the least cardinal with $2^\epsilon>\chi$.
Choose a smaller eligible level $\nu$ such that
\[
 |z|+\sigma+2^\chi<\nu<\mu,
\]
where $\mu$ is the level of $J$, with all required testing bounds
below $\nu$ and with
Theorem~\ref{iter:boundary-completion} available at $\nu$.
The further closure requirement on the level of $J$
ensures this choice. It is a finite refinement for this
fixed $n$.

Fix a well-ordered presentation $J=F_n(I)$. Construct
continuous subpresentations $J_i$ of size $\nu$ for
$i\le\epsilon$, beginning with $z$ in $J_0$.
At step $i$, the boundary of $J_i$ is a bounded parameter
set in $J$. By the assumed failure, a boundary
automorphism $\alpha_i$ of $J$ fixes it pointwise and
changes the value of the function at some tuple $u_i$
of length at most $\sigma$. Thus
\[
 D_J(z,u_i)\ne D_J(z,\alpha_i(u_i)).
 \tag{S}
\]
We choose $J_{i+1}$ containing both tests so that
$\alpha_i$ restricts to an automorphism of its boundary.

Here is the support closure for this choice. Begin with
the positions for $J_i$, $z$, $u_i$, and
$\alpha_i(u_i)$, padding to size $\nu$. For every element
of every proper projected subpresentation currently
generated, add positions for its image and its inverse
image under the corresponding face map of $\alpha_i$.
There are finitely many faces and $\nu$ elements in their
union, and each representation uses finitely many
positions. Repeat this operation countably many times.
The resulting position set still has size $\nu$.
Every element generated at its union appeared at a
finite stage, so both its image and inverse image are
generated in the union. Hence all restricted face maps
are onto and remain compatible with the specified
inclusions. They give the asserted boundary
automorphism of $J_{i+1}$. At a limit use the union of
the position sets. Mixed diagram preservation transfers
$(\mathrm S)$ to $J_{i+1}$, since $|z|+\sigma<\nu$.

Now build a binary tree of boundary embeddings
\[
 q_s:\partial J_i\longrightarrow\partial A,
 \qquad s\in{}^i2,
\]
where $A$ is an increasing accumulator of size $\nu$.
Start with $A=J_0$ and the identity map on $\partial J_0$.
At a successor node first apply completion to $q_s$
on $J_i$ and then amalgamate $J_{i+1}$ over that full
copy of $J_i$. This yields an arrow
$G_s:J_{i+1}\to A'$ whose boundary extends $q_s$.
Define the two next boundary embeddings by
\[
 q_{s0}=\partial G_s,
 \qquad
 q_{s1}=(\partial G_s)\alpha_i^{-1}.
\]
They restrict to the same old $q_s$, because
$\alpha_i$ fixes $\partial J_i$ pointwise. The common
parameter test $G_s(u_i)$ is the image of $u_i$ on
branch $s0$ and of $\alpha_i(u_i)$ on branch $s1$.
Only these boundary embeddings are retained for later
steps; the full map $G_s$ is temporary.

At limits the continuity of the $J_i$ and the
common-upper-bound clauses give the boundary maps into
the union accumulator. At height $\epsilon$, complete
the boundary maps to full maps of $J_\epsilon$, one
leaf at a time. The full maps at the leaves carry $z$.
If two leaves first separate at $s\in{}^i2$, their
images of $z$ have the two distinct diagrams from
$(\mathrm S)$ over the one retained test $G_s(u_i)$.
This follows from preservation by their final full maps
and from coherence of their boundary maps with the
chosen maps at the split.

There are at most $\chi$ nonterminal nodes, hence at
most $\chi$ parameters in all the splitting tests.
There are $2^\epsilon>\chi$ distinguishable leaf
images of $z$. All leaf completions and all unions
remain at level $\nu$, since $2^\chi<\nu$.
Lemma~\ref{calc:parameter-orbits} again contradicts
this number of orbits over the tests. Thus the bounded
fixing list $a$ exists.
\end{proof}

\begin{theorem}[Exact boundary lifting]
\label{iter:exact-boundary}
For each finite $n$, at further sufficiently closed
levels every boundary isomorphism
$h:\partial J\overset{\cong}{\longrightarrow}\partial K$
between $\mathcal C_n$-objects at the same level
extends to an isomorphism $J\cong K$ onto the given
target. All the prescribed marks and bounded named
data are preserved.
\end{theorem}

\begin{proof}
We build a back-and-forth sequence of matches
$a\mapsto a'$ of bounded tuples, maintaining
\[
 D_J(a,u)=D_K(a',h(u))
 \quad\text{for every bounded boundary tuple }u.
 \tag{BF}
\]
For the empty tuple this is
Lemma~\ref{iter:boundary-diagrams}. Suppose that a
match satisfying $(\mathrm{BF})$ has been chosen, and let $z$
be any bounded enlargement of $a$ in $J$.

Take a bounded infinite $\sigma$ at least the uniform
testing bound for a tuple of length $|z|$ over the
natural boundary layout of a $\mathcal C_n$-object.
Choose fixing data $c$ for $z,\sigma$ in $\partial J$
by Lemma~\ref{iter:boundary-fixed-parameters}.
The maintained equality with the particular boundary
list $c$ gives
\[
 D_J(a,c)=D_K(a',h(c)).
\]
Working-package homogeneity supplies a full
isomorphism $H:J\cong K$ carrying $a$ to $a'$ and
$c$ to $h(c)$. Put
$\alpha=(\partial H)^{-1}h$. This is a boundary
automorphism of $J$ fixing $c$ pointwise. Therefore,
for every boundary tuple $u$ of length at most
$\sigma$,
\[
 \begin{split}
 D_K(H(z),h(u))
 &=D_J(z,\alpha(u))\\
 &=D_J(z,u).
 \end{split}
 \tag{BT}
\]

We justify explicitly the passage from these tests
to all bounded boundary tests. By
Theorem~\ref{iter:boundary-completion}, embed $J$
into a same-level extension $K^*$ of $K$ by an arrow
$j$ extending $h$. Both $j(z)$ and $H(z)$ are now
tuples in $K^*$, and $(\mathrm{BT})$ says that their diagrams
agree over boundary tests from $K$ of size at most
$\sigma$. The boundary of $K$, in a well-ordered
chart of $K$, is a natural collection of term sorts.
The types of these two tuples over $K$ have order
placements by the one-dimensional calculus.
Proposition~\ref{calc:testing} restricted to this natural
boundary layout therefore gives equality over every
bounded boundary test. Preservation by $j$ then
shows
\[
 D_K(H(z),h(u))=D_J(z,u)
\]
for all bounded $u$. Thus $z\mapsto H(z)$ extends
the previous match and satisfies $(\mathrm{BF})$.

Apply this extension step alternately to $J,K$ with
$h$, and to $K,J$ with $h^{-1}$. Their level has
countable cofinality, so countably many bounded
enlargements can exhaust both systems. In each
vertex choose increasing smaller strong submodels
whose enumerations are included in these tuples.
The union assignment is onto in every component
and is an arrow by the working-package map
criterion; the reverse construction gives its
inverse arrow. If an entry lies on the boundary,
its equality test with that boundary entry in
$(\mathrm{BF})$ forces its image to be exactly the image
under $h$. The resulting isomorphism therefore
extends all of $h$, as required.
\end{proof}
